\newcommand{\DOCMATIERE}{Chimie}\newcommand{\DOCNIVEAU}{Tle C D E}\newcommand{\DOCMODULE}{Module 2 · Acides et bases}\newcommand{\DOCLECON}{09}\newcommand{\DOCTITRE}{Réactions acide/base et dosages}
% OnBuch+ — préambule « cours signés » ENRICHI (compilé avec tectonic / XeLaTeX)
% Système visuel « Pop » d'OnBuch+ : fond crème (jamais blanc), bordures encre
% épaisses, ombres « dures » décalées, titres Archivo Black en capitales.
% Version enrichie pour les cours de chimie : chimie (mhchem, chemfig),
% graphes (pgfplots), boîtes pédagogiques supplémentaires (définition,
% propriété, méthode, attention, le savais-tu, exercice résolu, exercice,
% TP, bilan), page de garde refaite et signature OnBuch+ ancrée en bas.
%
% Macros attendues dans main.tex AVANT \input{preamble} :
%   \DOCMATIERE  \DOCNIVEAU  \DOCMODULE  \DOCLECON (numéro)  \DOCTITRE
\documentclass[11pt]{article}

\usepackage{fontspec}
\usepackage[french]{babel}
\usepackage[a4paper,margin=2cm,top=3.4cm,bottom=2.8cm,headheight=1.4cm,footskip=1.3cm]{geometry}
\usepackage{amsmath,amssymb,amsfonts}
\usepackage[table]{xcolor} % [table] : fournit aussi \rowcolors (alternance de lignes)
\usepackage{pagecolor}
\usepackage{tikz}
\usetikzlibrary{arrows.meta,positioning,calc,shapes.geometric,shapes.misc,shapes.arrows,decorations.pathmorphing,decorations.markings,patterns,shadows,mindmap,trees,fit,backgrounds,matrix,intersections}
\usepackage{pgfplots}
\pgfplotsset{compat=1.18}
\usepackage[version=4]{mhchem}
\usepackage{chemfig}
\usepackage{enumitem}
\usepackage{fancyhdr}
% [most] charge toutes les bibliothèques tcolorbox (listings, minted, raster,
% documentation...) et gonfle inutilement la mémoire TeX sur un document long
% (~300 boîtes) : on ne charge que ce qui est réellement utilisé.
\usepackage[skins,breakable]{tcolorbox}
\usepackage{graphicx}
\usepackage{colortbl}
\usepackage{booktabs}
\usepackage{tabularx}
\usepackage{array}
\usepackage{multicol}
\usepackage{siunitx}
% parse-numbers=false : accepte aussi \SI{2,5 \times 10^{-3}}{...}
\sisetup{locale=FR,output-decimal-marker={,},group-minimum-digits=4,parse-numbers=false}
\DeclareSIUnit\ohm{\ensuremath{\symup{\Omega}}} % U+2126 absent de la police
\usepackage{qrcode}
% \degree : commande fréquemment utilisée par les modèles pour un angle ou une
% température en dehors de \SI{}{} (non fournie par les paquets chargés).
\providecommand{\degree}{^{\circ}}
\usepackage{lastpage}
\usepackage{hyperref}
\hypersetup{hidelinks}

% ============================================================
% Palette « Pop » OnBuch+ — mêmes hex que la classe P (plus_ui.dart)
% ============================================================
\definecolor{poporange}{HTML}{F59321}
\definecolor{popdark}{HTML}{E07A0C}
\definecolor{popcream}{HTML}{FFF6E8}
\definecolor{popink}{HTML}{1C1714}
\definecolor{poppurple}{HTML}{7A5AE0}
\definecolor{popgreen}{HTML}{2E9E6A}
\definecolor{poppink}{HTML}{E0457B}
\definecolor{popblue}{HTML}{2F7DE1}
\definecolor{popgold}{HTML}{C98A12}
\definecolor{popmuted}{HTML}{8A7F76}
\definecolor{popline}{HTML}{EADFD2}
\definecolor{popgray}{HTML}{8A7F76}
\colorlet{popgrey}{popgray}
% Alias tolérants (le modèle nomme parfois les couleurs différemment)
\colorlet{popwhite}{white}
\colorlet{popblack}{popink}
\colorlet{popred}{poppink}
\colorlet{popyellow}{popgold}
% Teintes claires pour fonds de tableaux / zones
\colorlet{popblueL}{popblue!10!white}
\colorlet{poporangeL}{poporange!14!white}
\colorlet{popgreenL}{popgreen!12!white}
\colorlet{poppurpleL}{poppurple!12!white}
\colorlet{poppinkL}{poppink!12!white}
\colorlet{popgoldL}{popgold!14!white}

\pagecolor{popcream}

% ============================================================
% Typographie
% ============================================================
\defaultfontfeatures{Ligatures=TeX}
\setmainfont{PlusJakartaSans-Medium.ttf}[Path=../../fonts/,BoldFont=PlusJakartaSans-Bold.ttf]
% Maths en sans-serif (Fira Math) avec chiffres et lettres droites de Plus
% Jakarta, pour que les formules se fondent dans le texte.
\usepackage[math-style=ISO,bold-style=ISO]{unicode-math}
\setmathfont{FiraMath-Regular.otf}[Path=../../fonts/]
\setmathfont{PlusJakartaSans-Medium.ttf}[Path=../../fonts/,range={up/num,up/{Latin,latin},\mathup}]
\usepackage{icomma} % virgule décimale sans espace en mode math
% Caractères absents de la police de texte (sinon carré vide) : rendus par
% leur équivalent mathématique, en texte comme en formule.
\usepackage{newunicodechar}
\newunicodechar{α}{\ensuremath{\alpha}}
\newunicodechar{Α}{\ensuremath{\alpha}}
\newunicodechar{β}{\ensuremath{\beta}}
\newunicodechar{δ}{\ensuremath{\delta}}
\newunicodechar{✓}{\popcheck{popgreen}}
\newfontfamily\popdisplay{ArchivoBlack-Regular.ttf}[Path=../../fonts/]
\newfontfamily\poplogo{SpaceGrotesk-Bold.ttf}[Path=../../fonts/]
\newfontfamily\popsemi{PlusJakartaSans-SemiBold.ttf}[Path=../../fonts/]
\newfontfamily\popxbold{PlusJakartaSans-ExtraBold.ttf}[Path=../../fonts/]
\newfontfamily\popxboldls{PlusJakartaSans-ExtraBold.ttf}[Path=../../fonts/,LetterSpace=3.0]

\setlist[enumerate]{leftmargin=1.7em,itemsep=5pt,topsep=4pt}
\setlist[itemize]{leftmargin=1.7em,itemsep=4pt,topsep=4pt,label={\color{poporange}\textbullet}}
\setlength{\parindent}{0pt}
\setlength{\parskip}{5pt}
\renewcommand{\arraystretch}{1.25}

% chemfig : liaisons un peu plus épaisses, encre Pop
\setchemfig{atom sep=2em,bond style={line width=0.9pt,popink},cram width=3pt}

% pgfplots : style Pop par défaut
\pgfplotsset{
  every axis/.append style={axis line style={popink,line width=1pt},tick style={popink},
    grid style={popline,line width=0.5pt},label style={font=\small},tick label style={font=\footnotesize}},
  popcurve/.style={poporange,line width=1.8pt,no markers,smooth},
}

% ============================================================
% Pill / squiggle
% ============================================================
\newcommand{\poppill}[3]{%
  \tikz[baseline=-0.55ex]\node[draw=popink,line width=1.2pt,rounded corners=2.6pt,%
    fill=#2,inner xsep=6.5pt,inner ysep=3pt]{%
    \popxboldls\fontsize{7}{7}\selectfont\color{#3}\MakeUppercase{#1}};%
}
\newcommand{\popsquiggle}[1]{%
  \tikz[baseline=-0.4ex]{
    \draw[#1,line width=2.6pt,line cap=round]
      plot [smooth] coordinates {(0,0.09) (0.34,0.01) (0.68,0.21) (1.02,0.01) (1.36,0.10)};
  }%
}
% Mot-clé surligné (marqueur orange sous le texte)
\newcommand{\cle}[1]{{\popxbold\color{popdark}#1}}

% ============================================================
% En-tête / pied
% ============================================================
\pagestyle{fancy}
\fancyhf{}
\renewcommand{\headrulewidth}{0pt}
\renewcommand{\footrulewidth}{0pt}
\lhead{\raisebox{-0.15ex}{\poplogo\fontsize{13}{13}\selectfont\color{popink}OnBuch\color{poporange}+}}
\rhead{\poppill{\DOCMATIERE\ \textbullet\ \DOCNIVEAU\ \textbullet\ Leçon \DOCLECON}{white}{popink}}
\lfoot{\popxbold\fontsize{7.6}{7.6}\selectfont\color{popmuted}\MakeUppercase{Cours signé OnBuch+}\ \textbullet\ {\popsemi\DOCTITRE}}
\rfoot{\tikz[baseline=-0.6ex]\node[rounded corners=8.5pt,fill=popink,minimum height=17pt,minimum width=17pt,inner xsep=5pt,inner sep=0]{\popxbold\fontsize{8}{8}\selectfont\color{popcream}\thepage\,/\,\pageref*{LastPage}};}

% ============================================================
% Titres
% ============================================================
\newcommand{\chaptitre}[2]{%
  \begin{tcolorbox}[enhanced,colback=poporange,colframe=popink,boxrule=2.6pt,arc=11pt,
      drop shadow=popink,left=15pt,right=15pt,top=12pt,bottom=11pt,before skip=3pt,after skip=18pt]
    \poppill{#2}{popink}{popcream}\par\vspace{9pt}
    {\raggedright\hyphenpenalty=10000\exhyphenpenalty=10000\popdisplay\fontsize{22}{24}\selectfont\color{popcream}\MakeUppercase{#1}\par}
  \end{tcolorbox}
}
% Section numérotée : pastille orange avec le numéro + titre Archivo
\newcounter{popsec}
\newcommand{\coursec}[1]{%
  \stepcounter{popsec}\setcounter{popsub}{0}%
  \par\vspace{16pt}\noindent
  \begin{minipage}{\linewidth}
  \tikz[baseline=-0.7ex]\node[draw=popink,line width=1.6pt,fill=poporange,rounded corners=4pt,minimum size=22pt,inner sep=0]{\popdisplay\fontsize{12}{12}\selectfont\color{popcream}\thepopsec};%
  \hspace{8pt}{\popdisplay\fontsize{15}{17}\selectfont\color{popink}\MakeUppercase{#1}}\par
  \vspace{4pt}\noindent{\color{poporange}\rule{36pt}{3.6pt}}
  \end{minipage}\par\nopagebreak\vspace{8pt}
}
\newcounter{popsub}
\newcommand{\courssub}[1]{%
  \stepcounter{popsub}%
  \par\vspace{9pt}\noindent{\popxbold\fontsize{11.5}{13}\selectfont\color{popink}{\color{poporange}\thepopsec.\thepopsub}\hspace{6pt}#1}\par\nopagebreak\vspace{4pt}
}

% ============================================================
% Boîtes pédagogiques — même recette : fond blanc, bordure épaisse colorée,
% ombre dure encre, badge plein. Titre optionnel [#1].
% ============================================================
\tcbset{popbase/.style={breakable,enhanced,colback=white,boxrule=2.3pt,arc=9pt,
  shadow={3pt}{-3pt}{0pt}{popink},left=14pt,right=14pt,top=11pt,bottom=11pt,before skip=11pt,after skip=15pt,
  fontupper=\popsemi\fontsize{10.6}{14.5}\selectfont\color{popink},
  fontlower=\popsemi\fontsize{10.6}{14.5}\selectfont\color{popink}}}
% Coche dessinée (le glyphe ✓ n'existe pas dans les polices du document)
\newcommand{\popcheck}[1]{\tikz[baseline=-0.5ex]\draw[#1,line width=1.6pt,line cap=round,line join=round](-0.13,0.0)--(-0.04,-0.09)--(0.13,0.1);}
\newcommand{\popboxhead}[3]{\poppill{#1}{#2}{white}\ifx\relax#3\relax\else\hspace{6pt}{\popxbold\fontsize{10.6}{12}\selectfont\color{popink}#3}\fi\par\vspace{7pt}}

\newtcolorbox{aretenir}[1][]{popbase,colframe=popgreen,before upper={\popboxhead{À retenir}{popgreen}{#1}}}
\newtcolorbox{exemplebox}[1][]{popbase,colframe=poporange,before upper={\popboxhead{Exemple}{poporange}{#1}}}
\newtcolorbox{definition}[1][]{popbase,colframe=popblue,before upper={\popboxhead{Définition}{popblue}{#1}}}
\newtcolorbox{propriete}[1][]{popbase,colframe=poppurple,before upper={\popboxhead{Propriété}{poppurple}{#1}}}
\newtcolorbox{methode}[1][]{popbase,colframe=popgold,before upper={\popboxhead{Méthode}{popgold}{#1}}}
\newtcolorbox{attention}[1][]{popbase,colframe=poppink,before upper={\popboxhead{Attention}{poppink}{#1}}}
\newtcolorbox{experience}[1][]{popbase,colframe=popblue,colback=popblueL,before upper={\popboxhead{Expérience}{popblue}{#1}}}
% « Le savais-tu ? » : carte violette pleine (contexte, culture, Cameroun)
\newtcolorbox{savaistu}[1][]{popbase,colback=poppurple,colframe=popink,
  fontupper=\popsemi\fontsize{10.4}{14.2}\selectfont\color{white},
  before upper={\poppill{Le savais-tu\,?}{white}{poppurple}\ifx\relax#1\relax\else\hspace{6pt}{\popxbold\color{white}#1}\fi\par\vspace{7pt}}}
% Exercice résolu : énoncé en haut, \tcblower puis la solution sur fond crème
\newcounter{popexo}\newcounter{popexr}
\newtcolorbox{exoresolu}[1][]{popbase,colframe=poporange,colbacklower=popcream,
  segmentation style={popink,line width=1.2pt,dashed},
  before upper={\stepcounter{popexr}\popboxhead{Exercice résolu \thepopexr}{poporange}{#1}},
  before lower={\poppill{Solution}{popink}{popcream}\par\vspace{6pt}}}
% Exercice (non résolu en place) — la correction est dans la partie Corrigés
\newtcolorbox{exercice}[1][]{popbase,colframe=popink,
  before upper={\stepcounter{popexo}\popboxhead{Exercice \thepopexo}{popink}{#1}}}
% Corrigé
\newtcolorbox{corrige}[1][]{popbase,colframe=popgreen,colback=popgreenL,
  before upper={\popboxhead{Corrigé}{popgreen}{#1}}}
% Objectifs / prérequis / bilan
\newtcolorbox{objectifs}{popbase,colframe=popink,colback=poporangeL,
  before upper={\popboxhead{Objectifs de la leçon}{popink}{}}}
\newtcolorbox{prerequis}{popbase,colframe=popink,colback=popblueL,
  before upper={\popboxhead{Prérequis}{popblue}{}}}
\newtcolorbox{bilan}{popbase,colframe=popink,colback=white,boxrule=2.8pt,
  before upper={\popboxhead{Fiche bilan}{poporange}{L'essentiel à connaître pour l'examen}}}
% Figure encadrée avec légende
\newtcolorbox{popfigure}[1][]{enhanced,breakable=false,colback=white,colframe=popline,boxrule=1.4pt,arc=8pt,
  left=10pt,right=10pt,top=10pt,bottom=8pt,before skip=10pt,after skip=12pt,halign=center,
  lower separated=false,halign lower=center,
  fontlower=\popsemi\fontsize{9}{11.5}\selectfont\color{popmuted},
  #1}
\newcommand{\legende}[1]{\par\vspace{6pt}{\popsemi\fontsize{9}{11.5}\selectfont\color{popmuted}#1}}

% ============================================================
% Signature OnBuch+
% ============================================================
\newcommand{\poplogoinline}[1]{{\poplogo\fontsize{#1}{#1}\selectfont\color{popink}OnBuch\color{poporange}+}}
\newcommand{\signatureonbuch}{%
  \begin{tcolorbox}[enhanced,colback=popink,colframe=popink,boxrule=2.2pt,arc=10pt,
      drop shadow=poporange,left=14pt,right=14pt,top=10pt,bottom=10pt,before skip=0pt,after skip=0pt]
    \tikz[baseline=-0.7ex]\node[circle,fill=popgreen,draw=popcream,line width=1.3pt,minimum size=22pt,inner sep=0]{\popcheck{white}};%
    \hspace{9pt}\begin{minipage}[c]{0.62\linewidth}
      {\popxbold\fontsize{12}{13}\selectfont\color{popcream}Signé\ \textbullet\ {\poplogo OnBuch\color{poporange}+}}\par\vspace{2pt}
      {\raggedright\popsemi\fontsize{8}{10.5}\selectfont\color{popline}\MakeUppercase{Équipe pédagogique OnBuch+}\\\MakeUppercase{Conforme au programme officiel MINESEC}\par}
    \end{minipage}\hfill
    {\poplogo\fontsize{22}{22}\selectfont\color{popcream}OnBuch\color{poporange}+}
  \end{tcolorbox}%
}

% ============================================================
% Page de garde
% #1 titre · #2 matière · #3 niveau · #4 module · #5 n° leçon · #6 durée ·
% #7 description / objectifs (texte ou itemize)
% ============================================================
\newcommand{\pagegarde}[7]{%
  \thispagestyle{empty}
  \newgeometry{a4paper,margin=1.7cm,top=1.6cm,bottom=1.4cm}%
  \begin{tikzpicture}[remember picture,overlay]
    \fill[poporange!18] ([xshift=-3.2cm,yshift=-3.6cm]current page.north east) circle (4.6cm);
    \fill[poppurple!14] ([xshift=2.4cm,yshift=7.6cm]current page.south west) circle (3.8cm);
  \end{tikzpicture}%
  {\poplogo\fontsize{30}{30}\selectfont\color{popink}OnBuch\color{poporange}+}\hfill
  \raisebox{0.6ex}{\poppill{Cours signé \textbullet\ Premium}{popink}{popcream}}\par
  \vspace{1pt}\popsquiggle{poppurple}\par
  \vspace{8pt}
  \begin{tcolorbox}[enhanced,colback=poporange,colframe=popink,boxrule=2.8pt,arc=13pt,
      drop shadow=popink,left=18pt,right=18pt,top=12pt,bottom=13pt,after skip=10pt]
    \poppill{#2 \textbullet\ #3}{popink}{popcream}\hspace{5pt}\poppill{#4}{white}{popink}\par\vspace{8pt}
    {\popdisplay\fontsize{14}{14}\selectfont\color{popink}LEÇON #5}\par\vspace{4pt}
    {\raggedright\hyphenpenalty=10000\exhyphenpenalty=10000\popdisplay\fontsize{25}{27}\selectfont\color{popcream}\MakeUppercase{#1}\par}
    \vspace{8pt}
    {\popxbold\fontsize{9.5}{9.5}\selectfont\color{popink}\MakeUppercase{Durée conseillée : #6}}
  \end{tcolorbox}
  \begin{tcolorbox}[enhanced,colback=white,colframe=popink,boxrule=2.2pt,arc=10pt,
      drop shadow=popink,left=16pt,right=16pt,top=10pt,bottom=10pt,after skip=8pt]
    \poppill{Dans cette leçon}{poppurple}{white}\par\vspace{5pt}
    \popsemi\fontsize{9.3}{12.6}\selectfont\color{popink}#7
  \end{tcolorbox}
  \noindent\begin{minipage}[t]{0.487\linewidth}
    \begin{tcolorbox}[enhanced,colback=popgreen,colframe=popink,boxrule=2.2pt,arc=10pt,
        drop shadow=popink,left=11pt,right=11pt,top=10pt,bottom=10pt,height=3.1cm,valign=center]
      \tcbox[colback=white,colframe=popink,boxrule=1.3pt,arc=5pt,boxsep=0pt,nobeforeafter,
        left=3pt,right=3pt,top=3pt,bottom=3pt,valign=center]{\qrcode[height=1.9cm]{https://play.google.com/store/apps/details?id=com.luvvix.onbuch&hl=fr}}%
      \hspace{8pt}\begin{minipage}[c]{0.5\linewidth}
        {\popdisplay\fontsize{11}{12.5}\selectfont\color{white}\MakeUppercase{Télécharge OnBuch}}\par\vspace{4pt}
        {\raggedright\popsemi\fontsize{8.4}{11}\selectfont\color{white}Gratuit \textbullet\ Google Play\\Tuteur IA, annales, concours, résultats\par}
      \end{minipage}
    \end{tcolorbox}
  \end{minipage}\hfill
  \begin{minipage}[t]{0.487\linewidth}
    \begin{tcolorbox}[enhanced,colback=poppurple,colframe=popink,boxrule=2.2pt,arc=10pt,
        drop shadow=popink,left=11pt,right=11pt,top=10pt,bottom=10pt,height=3.1cm,valign=center]
      \tikz[baseline=-0.7ex]\node[circle,fill=white,draw=popink,line width=1.3pt,minimum size=1.6cm,inner sep=0]{\popdisplay\fontsize{24}{24}\selectfont\color{poppurple}+};%
      \hspace{8pt}\begin{minipage}[c]{0.6\linewidth}
        {\popdisplay\fontsize{11}{12.5}\selectfont\color{white}\MakeUppercase{Passe à OnBuch+}}\par\vspace{4pt}
        {\raggedright\popsemi\fontsize{8.4}{11}\selectfont\color{white}Tous les cours signés, Tuteur IA illimité, fiches PDF — onglet OnBuch+ de l'app\par}
      \end{minipage}
    \end{tcolorbox}
  \end{minipage}
  \vspace{8pt}
  \popcontenu
  \vfill
  \signatureonbuch
  \restoregeometry
  \newpage
}

% Rangée « Ce cours contient » (page de garde)
\newcommand{\poptile}[3]{% #1 couleur #2 gros texte #3 légende
  \begin{tcolorbox}[enhanced,colback=#1,colframe=popink,boxrule=2pt,arc=9pt,shadow={2.5pt}{-2.5pt}{0pt}{popink},
      left=6pt,right=6pt,top=9pt,bottom=9pt,halign=center,valign=center,height=1.75cm,width=\linewidth,nobeforeafter]
    {\popdisplay\fontsize{12.5}{13}\selectfont\color{white}\MakeUppercase{#2}}\par\vspace{3pt}
    {\centering\popsemi\fontsize{7.8}{9.5}\selectfont\color{white}#3\par}
  \end{tcolorbox}}
\newcommand{\popcontenu}{%
  \par\noindent{\popxboldls\fontsize{7.5}{7.5}\selectfont\color{popmuted}CE COURS SIGNÉ CONTIENT}\par\vspace{7pt}
  \noindent\begin{minipage}[t]{0.235\linewidth}\poptile{popblue}{Cours}{complet, illustré, pas à pas}\end{minipage}\hfill
  \begin{minipage}[t]{0.235\linewidth}\poptile{poporange}{Méthodes}{et exercices résolus}\end{minipage}\hfill
  \begin{minipage}[t]{0.235\linewidth}\poptile{poppink}{Exercices}{type examen + corrigés}\end{minipage}\hfill
  \begin{minipage}[t]{0.235\linewidth}\poptile{popgreen}{Fiche bilan}{l'essentiel à retenir}\end{minipage}\par}

% Sommaire
\newtcolorbox{sommaire}{enhanced,colback=white,colframe=popink,boxrule=2.2pt,arc=10pt,
  drop shadow=popink,left=18pt,right=18pt,top=13pt,bottom=13pt,before skip=6pt,after skip=20pt,
  before upper={\poppill{Sommaire}{poppurple}{white}\par\vspace{9pt}\popsemi\fontsize{10.6}{14.5}\selectfont\color{popink}}}

% Signature de fin de cours — poussée tout en bas de la dernière page
\newcommand{\findecours}{%
  \par\vfill\nopagebreak
  \begin{center}\popsquiggle{poporange}\end{center}\vspace{4pt}
  \signatureonbuch
}
% Glyphes absents de Fira Math : redéfinis à partir de glyphes disponibles
\AtBeginDocument{%
  \def\setminus{\mathbin{\backslash}}\let\smallsetminus\setminus
  \def\vdots{\mathord{\vbox{\baselineskip3.2pt\lineskiplimit0pt\kern4pt\hbox{.}\hbox{.}\hbox{.}}}}%
  \def\ddots{\mathinner{\mkern1mu\raise6.5pt\hbox{.}\mkern2mu\raise3.6pt\hbox{.}\mkern2mu\raise0.7pt\hbox{.}\mkern1mu}}%
  \def\triangle{\mathord{\scalebox{0.9}{$\symup{\Delta}$}}}%
  \def\nabla{\mathord{\raisebox{\depth}{\scalebox{1}[-1]{$\symup{\Delta}$}}}}%
  \def\checkmark{\popcheck{popdark}}%
  \def\star{\mathbin{\ast}}\def\bigstar{\mathord{\ast}}%
  \def\diamond{\mathbin{\scalebox{0.6}{$\Diamond$}}}%
  \def\bigcup{\mathop{\mathchoice{\raisebox{-0.3em}{\scalebox{1.7}{$\cup$}}}{\scalebox{1.25}{$\cup$}}{\cup}{\cup}}\displaylimits}%
  \def\bigcap{\mathop{\mathchoice{\raisebox{-0.3em}{\scalebox{1.7}{$\cap$}}}{\scalebox{1.25}{$\cap$}}{\cap}{\cap}}\displaylimits}%
  \def\lesseqqgtr{\mathrel{\lessgtr}}\def\bigoplus{\mathop{\oplus}\displaylimits}%
}
\providecommand{\ssi}{si et seulement si} % abréviation fréquente

\begin{document}
\begin{tcolorbox}[enhanced,colback=popink,colframe=popink,arc=10pt,boxrule=0pt,left=12pt,right=12pt,top=9pt,bottom=9pt]
{\color{white}\popxbold\small FICHE DE TD 3/3 \textbullet\ DIFFICILE \hfill Chimie \textbullet\ Tle C D E\par}
\vspace{4pt}{\color{white}\popxbold\Large Réactions acide/base et dosages\par}
\vspace{3pt}{\color{white}\small Module 2 · Acides et bases\par}
\end{tcolorbox}
\vspace{4pt}
Ce TD de synthèse mobilise l'ensemble des notions du module : étude critique de courbes, choix d'indicateurs, calculs de concentrations, hydrolyse à l'équivalence et vérification d'étiquette.

\begin{exercice}[Analyse comparative de trois dosages pH-métriques]
On réalise trois dosages pH-métriques distincts en versant une solution basique de concentration connue $C_B = \SI{0,100}{mol.L^{-1}}$ dans \SI{20,0}{mL} de trois solutions acides différentes (A, B, C) de concentration molaire identique $C_A = \SI{0,100}{mol.L^{-1}}$.
Les systèmes acides sont : \ce{HCl} (acide fort), \ce{CH3COOH} (acide éthanoïque, $pK_a = 4,80$) et \ce{NH4+} (acide conjugué de l'ammoniac, $pK_a = 9,20$).
Les courbes $\text{pH} = f(V_B)$ obtenues présentent les caractéristiques suivantes au point d'équivalence ($V_{BE} = \SI{20,0}{mL}$) :
\begin{itemize}
    \item Courbe 1 : pH = 7,0 ; saut de pH important (environ 3 à 11).
    \item Courbe 2 : pH = 8,7 ; saut de pH modéré (environ 7 à 10).
    \item Courbe 3 : pH = 5,1 ; saut de pH réduit (environ 4 à 7).
\end{itemize}
\begin{enumerate}
    \item Associer chaque courbe au système acide correspondant. Justifier la valeur du pH à l'équivalence par l'écriture de la réaction d'hydrolyse prépondérante et le calcul théorique approché (hypothèse $C_{\text{sel}} \approx \SI{0,050}{mol.L^{-1}}$).
    \item Pour le dosage de l'acide éthanoïque (Courbe 2), déterminer graphiquement (méthode des tangentes parallèles) le volume d'équivalence $V_{BE}$ et le volume de demi-équivalence $V_{1/2}$. Vérifier la relation $\text{pH}_{1/2} = pK_a$.
    \item Choisir, parmi l'hélianthine (virage 3,1--4,6), le BBT (virage 6,0--7,6) et la phénolphtaléine (virage 8,2--10,0), l'indicateur coloré adapté pour \textbf{chaque} dosage. Justifier rigoureusement chaque choix en lien avec la zone de virage et le saut de pH.
    \item On souhaite doser un mélange des trois acides (même volume, même concentration) par la soude \SI{0,100}{mol.L^{-1}}. Prédire l'allure de la courbe $\text{pH} = f(V_B)$ : nombre de sauts, volumes d'équivalence successifs, pH aux équivalences. Préciser si un dosage distinct de chaque espèce est possible par pH-métrie.
\end{enumerate}
\end{exercice}

\begin{exercice}[Vérification de l'étiquetage d'un vinaigre commercial : dosage par dilution]
Un vinaigre commercial porte la mention : « Acidité : 8,0 \% en masse d'acide éthanoïque » (masse volumique $\rho = \SI{1,01}{g.mL^{-1}}$).
On prélève \SI{10,0}{mL} de ce vinaigre (solution mère $S_0$) que l'on introduit dans une fiche jaugée de \SI{100,0}{mL} (solution fille $S_1$). On titre \SI{20,0}{mL} de $S_1$ par une solution de soude \ce{NaOH} de concentration $C_B = \SI{0,100}{mol.L^{-1}}$.
Le suivi pH-métrique donne les points suivants au voisinage de l'équivalence :
\begin{center}
\begin{tabular}{c|cccccccc}
$V_B$ (mL) & 14,0 & 15,0 & 15,5 & 15,8 & 16,0 & 16,2 & 16,5 & 17,0 \\
\hline
pH & 5,2 & 6,1 & 7,2 & 8,5 & 8,9 & 9,3 & 10,1 & 10,8
\end{tabular}
\end{center}
\begin{enumerate}
    \item Écrire l'équation-bilan de la réaction de dosage. Exprimer la relation liant les quantités de matière à l'équivalence.
    \item Tracer la courbe $\text{pH} = f(V_B)$ (échelle adaptée). Déterminer $V_{BE}$ par la méthode des tangentes parallèles (préciser la construction graphique). En déduire la concentration molaire $C_1$ de la solution $S_1$ et la concentration molaire $C_0$ de la solution mère $S_0$.
    \item Calculer la concentration massique de l'acide éthanoïque dans le vinaigre pur ($S_0$) en \SI{}{g.L^{-1}} et en \% massique. Conclure sur la conformité de l'étiquetage (tolérance admise $\pm 0,5\%$).
    \item Sur la courbe tracée, repérer le point de demi-équivalence. En déduire une valeur expérimentale du $pK_a$ de l'acide éthanoïque. Commenter la précision de cette méthode.
    \item Justifier le choix de la phénolphtaléine comme indicateur pour ce dosage colorimétrique. Que se passerait-il si l'on utilisait l'hélianthine ? Calculer l'erreur relative sur le volume équivalent si l'on stoppe au virage de l'hélianthine (pH $\approx 4,6$), en supposant que la courbe suit $\text{pH} = pK_a + \log(\frac{V_B}{V_{BE}-V_B})$ avant équivalence.
\end{enumerate}
\end{exercice}

\begin{exercice}[Dosage d'une solution d'ammoniac : méthode de la dérivée et hydrolyse]
On titre \SI{20,0}{mL} d'une solution d'ammoniac \ce{NH3(aq)} (base faible, $K_b = \SI{1,8e-5}{}$) par une solution d'acide chlorhydrique \ce{HCl} de concentration $C_A = \SI{0,100}{mol.L^{-1}}$.
Le tableau ci-dessous donne le pH en fonction du volume d'acide versé $V_A$ :
\begin{center}
\begin{tabular}{c|cccccccccccc}
$V_A$ (mL) & 0,0 & 5,0 & 10,0 & 15,0 & 18,0 & 19,0 & 19,5 & 20,0 & 20,5 & 21,0 & 22,0 & 25,0 \\
\hline
pH & 11,1 & 10,3 & 9,7 & 9,2 & 8,6 & 8,0 & 7,4 & 5,3 & 4,8 & 4,4 & 3,9 & 3,1
\end{tabular}
\end{center}
\begin{enumerate}
    \item Écrire l'équation-bilan de la réaction de dosage. Préciser la nature de la réaction (totale/partielle, rapide/lente, effet thermique).
    \item Tracer la courbe $\text{pH} = f(V_A)$. Déterminer $V_{AE}$ par la méthode des tangentes parallèles.
    \item La méthode de la dérivée consiste à tracer $\frac{\Delta \text{pH}}{\Delta V_A} = f(V_A)$ (ou $\frac{d\text{pH}}{dV_A}$). Expliquer pourquoi le maximum de cette courbe correspond à l'équivalence. Sans tracer la dérivée, estimer $V_{AE}$ en calculant les taux de variation moyens $\frac{\Delta \text{pH}}{\Delta V_A}$ sur les intervalles $[\SI{19,0}{mL};\SI{19,5}{mL}]$, $[\SI{19,5}{mL};\SI{20,0}{mL}]$, $[\SI{20,0}{mL};\SI{20,5}{mL}]$. Comparer les deux méthodes.
    \item Calculer la concentration molaire $C_B$ de la solution d'ammoniac.
    \item Calculer le pH théorique à l'équivalence. (Indication : l'espèce présente est \ce{NH4+}, acide faible de $K_a = K_e/K_b$ ; concentration du sel $\approx \SI{0,050}{mol.L^{-1}}$). Comparer à la valeur expérimentale lue sur la courbe.
    \item Parmi l'hélianthine, le BBT et la phénolphtaléine, quel(s) indicateur(s) permettent un dosage colorimétrique rigoureux ? Justifier en comparant les zones de virage au saut de pH expérimental et théorique.
    \item On souhaite préparer une solution tampon de pH = 9,20 à partir de cette solution d'ammoniac et d'acide chlorhydrique. Quel volume d'acide \SI{0,100}{mol.L^{-1}} faut-il ajouter à \SI{50,0}{mL} de solution d'ammoniac ? Justifier par le calcul.
\end{enumerate}
\end{exercice}

\begin{exercice}[Synthèse : Constante d'équivalence, constante d'hydrolyse et domaine d'indicateurs]
Soit le dosage d'un acide $\ce{AH}$ (concentration $C_A$, volume $V_A$) par une base $\ce{B}$ (concentration $C_B$, volume $V_B$). On note $K_A$ la constante d'acidité de $\ce{AH}$ et $K_B$ la constante de basicité de $\ce{B}$ ($K_A \times K_{A/B} = K_e = \SI{1,0e-14}{}$ à \SI{25}{\celsius}).
\begin{enumerate}
    \item Rappeler l'expression de la constante d'équivalence $K$ de la réaction $\ce{AH + B <=> A- + BH+}$ en fonction de $K_A$, $K_B$ et $K_e$. Discuter de la « totalité » de la réaction selon les couples acide/base mis en jeu (Fort/Fort, Faible/Fort, Fort/Faible, Faible/Faible).
    \item Cas \textbf{Acide faible / Base forte} ($K_A \ll 1, K_B \gg 1$). Montrer que le pH à l'équivalence s'écrit : $\text{pH}_E = \frac{1}{2}pK_e + \frac{1}{2}pK_A - \frac{1}{2}\log C_{\text{sel}}$ (hypothèses à préciser). Application numérique pour l'acide éthanoïque \SI{0,10}{mol.L^{-1}} titré par soude \SI{0,10}{mol.L^{-1}}.
    \item Cas \textbf{Acide fort / Base faible} ($K_A \gg 1, K_B \ll 1$). Établir l'expression symétrique du $\text{pH}_E$ en fonction de $pK_B$ et $C_{\text{sel}}$. Application pour l'ammoniac \SI{0,10}{mol.L^{-1}} titré par HCl \SI{0,10}{mol.L^{-1}}.
    \item On dispose d'indicateurs aux virages : I1 (3,0--4,6), I2 (6,0--7,6), I3 (8,2--10,0). Pour chacun des 4 cas de la question 1, indiquer quel(s) indicateur(s) conviennent. Justifier en analysant la position du pH à l'équivalence par rapport à la zone de virage et l'amplitude du saut de pH (liée à $K$).
    \item Problème inverse : On titre \SI{20,0}{mL} d'une solution d'acide inconnu (fort ou faible) par de la soude \SI{0,100}{mol.L^{-1}}. L'équivalence est trouvée à \SI{20,0}{mL}. Le pH à l'équivalence est 8,70. La courbe présente un point de demi-équivalence à \SI{10,0}{mL} pour un pH de 4,80.
    Déterminer la nature de l'acide (fort ou faible), sa concentration $C_A$, et son $pK_a$ s'il est faible. Quel indicateur choisir ?
\end{enumerate}
\end{exercice}
\end{document}
